\documentclass[11pt]{article}

\usepackage[T1]{fontenc}
\usepackage[utf8]{inputenc}
\usepackage[spanish,es-nodecimaldot]{babel}
\usepackage{amsmath,amssymb,mathtools}
\usepackage{enumitem}
\usepackage[a4paper,margin=2.5cm]{geometry}
\usepackage[hidelinks]{hyperref}
\usepackage{microtype}

% -----------------------------------------------------------------------------
% Notación
% -----------------------------------------------------------------------------
\newcommand{\C}{\mathbb{C}}
\newcommand{\N}{\mathbb{N}}
\newcommand{\BH}{\mathcal{B}(H)}
\newcommand{\spann}{\operatorname{span}}
\newcommand{\rank}{\operatorname{rank}}
\newcommand{\id}{\mathrm{id}}

% Entornos simples: deliberadamente sin cajas ni fondos de color.
\newenvironment{definicion}{\par\medskip\noindent\textbf{Definición.}\ }{\par\medskip}
\newenvironment{ejemplo}{\par\medskip\noindent\textbf{Ejemplo.}\ }{\par\medskip}
\newenvironment{observacion}{\par\medskip\noindent\textbf{Observación.}\ }{\par\medskip}
\newenvironment{ejercicio}{\par\medskip\noindent\textbf{Ejercicio.}\ }{\par\medskip}
\newenvironment{calculo}{\par\medskip\noindent\textbf{Cálculo.}\ }{\par\medskip}

\title{Sesión del 19 de octubre\\Ejemplos concretos de $C^*$-álgebras}
\author{}
\date{}

\begin{document}
\maketitle

\noindent
Estas notas son una transcripción depurada de las fotografías de la sesión. Se han
eliminado únicamente repeticiones fotográficas del mismo desarrollo; se conserva
la notación y el orden matemático de la clase. Cuando una igualdad de las fotos
presenta una ambigüedad de conjugación, se señala explícitamente en una nota de
transcripción al final.

\section{Convención para el producto interno}

Sea $H$ un espacio de Hilbert complejo y sea $(e_i)$ una base ortonormal de $H$.
La convención utilizada en la sesión es que el producto interno es \emph{lineal en
la primera variable}. En particular,
\begin{align*}
\langle \alpha h_1+\beta h_2,g\rangle
 &=\alpha\langle h_1,g\rangle+\beta\langle h_2,g\rangle,\\
\langle h,h\rangle&\ge 0,
\qquad
\langle h,h\rangle=0\iff h=0,\\
\langle g,h\rangle&=\overline{\langle h,g\rangle}.
\end{align*}
Por consiguiente, la segunda variable es conjugado-lineal:
\[
\langle h,\alpha g_1+\beta g_2\rangle
=
\overline{\alpha}\,\langle h,g_1\rangle
+
\overline{\beta}\,\langle h,g_2\rangle.
\]

\section{Operadores de rango uno en $\mathcal{B}(H)$}

\begin{definicion}
Para $i,j$ se define el operador
\[
 e_i\otimes e_j:H\longrightarrow H,
 \qquad
 (e_i\otimes e_j)(h):=\langle h,e_j\rangle e_i.
\]
En estas notas el símbolo $e_i\otimes e_j$ denota este operador de rango uno.
\end{definicion}

Su imagen es
\[
(e_i\otimes e_j)(H)=\C e_i,
\]
y, por tanto,
\[
\dim (e_i\otimes e_j)(H)=1.
\]
Así, $e_i\otimes e_j$ es un operador de rango uno.

Además, por Cauchy--Schwarz y porque $\|e_i\|=\|e_j\|=1$,
\begin{align*}
\|(e_i\otimes e_j)(h)\|
 &=\|\langle h,e_j\rangle e_i\|\\
 &=|\langle h,e_j\rangle|\,\|e_i\|\\
 &\le \|h\|\,\|e_j\|\,\|e_i\|\\
 &=\|h\|.
\end{align*}
En particular,
\[
\|e_i\otimes e_j\|\le 1,
\qquad
 e_i\otimes e_j\in\BH.
\]

\subsection{Adjunto de los operadores $e_i\otimes e_j$}

Escribamos
\[
T_{ij}:=e_i\otimes e_j.
\]
Entonces
\[
T_{ij}^*=T_{ji},
\qquad\text{es decir,}\qquad
(e_i\otimes e_j)^*=e_j\otimes e_i.
\]

En efecto, para $h,g\in H$,
\begin{align*}
\langle T_{ij}h,g\rangle
 &=\langle \langle h,e_j\rangle e_i,g\rangle\\
 &=\langle h,e_j\rangle\langle e_i,g\rangle\\
 &=\langle h,\langle g,e_i\rangle e_j\rangle\\
 &=\langle h,T_{ji}g\rangle.
\end{align*}
Por la definición del adjunto se obtiene $T_{ij}^*=T_{ji}$.

\subsection{Regla de multiplicación}

Para índices $i,j,r,s$,
\[
(e_i\otimes e_j)(e_r\otimes e_s)
=
\delta_{jr}\,e_i\otimes e_s.
\]
En efecto, para $h\in H$,
\begin{align*}
(e_i\otimes e_j)(e_r\otimes e_s)(h)
 &= (e_i\otimes e_j)\bigl(\langle h,e_s\rangle e_r\bigr)\\
 &=\langle h,e_s\rangle\langle e_r,e_j\rangle e_i\\
 &=\delta_{jr}\langle h,e_s\rangle e_i.
\end{align*}

Esta es la misma regla que satisfacen las unidades matriciales.

\subsection{Composición con un operador arbitrario}

Sea $T\in\BH$. En las fotografías se dejan como ejercicio las identidades
\[
T\circ(e_i\otimes e_j)=T(e_i)\otimes e_j
\]
y
\[
(e_i\otimes e_j)\circ T=e_i\otimes T^*(e_j).
\]

La segunda identidad se lee directamente de
\[
((e_i\otimes e_j)T)(h)
 =\langle Th,e_j\rangle e_i
 =\langle h,T^*e_j\rangle e_i.
\]

\section{Proyecciones de rango finito}

Para cada $i$,
\[
p_i:=e_i\otimes e_i
\]
es una proyección de $\BH$, pues
\[
p_i^*=p_i
\]
y, usando la regla de multiplicación,
\[
p_i^2=(e_i\otimes e_i)(e_i\otimes e_i)=e_i\otimes e_i=p_i.
\]
Además,
\[
p_i(H)=\C e_i,
\]
de modo que $p_i$ tiene rango uno.

\subsection{La esquina determinada por una proyección de rango uno}

Sea $p_i=e_i\otimes e_i$ y sea $T\in\BH$. Escribamos
\[
T(e_i)=\sum_r \lambda_r e_r.
\]
Entonces la esquina
\[
p_i\BH p_i
\]
consiste en múltiplos escalares de $p_i$. En particular,
\[
p_i\BH p_i=\C p_i\cong\C,
\qquad
\lambda p_i\longmapsto \lambda.
\]

Una forma directa de ver el coeficiente es
\begin{align*}
p_iTp_i(h)
 &=p_iT\bigl(\langle h,e_i\rangle e_i\bigr)\\
 &=\langle h,e_i\rangle p_iT(e_i)\\
 &=\langle h,e_i\rangle
   \left\langle \sum_r\lambda_r e_r,e_i\right\rangle e_i\\
 &=\lambda_i\langle h,e_i\rangle e_i\\
 &=\lambda_i p_i(h).
\end{align*}
Por tanto,
\[
p_iTp_i=\lambda_i p_i.
\]

\begin{ejercicio}
Si $A$ es una $C^*$-álgebra y $p\in A$ es una proyección, comprobar que
\[
pAp:=\{pap:a\in A\}
\]
es una $C^*$-subálgebra de $A$. En las fotografías esta subálgebra se denomina
\emph{corner} o \emph{esquina} determinada por $p$.
\end{ejercicio}

En particular, cuando $p$ es una proyección de rango uno en una representación
sobre un espacio de Hilbert,
\[
pAp\cong\C
\]
en la situación considerada en clase.

\subsection{Rango de una proyección}

En la sesión se usa la terminología
\[
p\text{ es de rango uno}
\quad\Longleftrightarrow\quad
\dim p(H)=1,
\]
y, más generalmente,
\[
p\text{ es de rango }n
\quad\Longleftrightarrow\quad
\dim p(H)=n.
\]

\begin{ejemplo}
Consideremos
\[
p=e_1\otimes e_1+e_2\otimes e_2.
\]
Usando la regla de multiplicación de las unidades matriciales,
\begin{align*}
p^2
 &= (e_1\otimes e_1+e_2\otimes e_2)^2\\
 &=e_1\otimes e_1+e_2\otimes e_2\\
 &=p.
\end{align*}
Además $p^*=p$, y
\[
p(H)=\C e_1+\C e_2\cong\C^2.
\]
Por tanto, $p$ es una proyección de rango $2$.
\end{ejemplo}

\section{La esquina $p\mathcal{B}(H)p$ para una proyección de rango dos}

Sea
\[
p=e_1\otimes e_1+e_2\otimes e_2.
\]
Para $T\in\BH$, escribimos las coordenadas de $T$ respecto de la base ortonormal
como
\[
T(e_r)=\sum_k\lambda_{kr}e_k.
\]
Entonces
\begin{align*}
(e_i\otimes e_j)T(e_r\otimes e_s)
 &= (e_i\otimes e_j)\bigl(T(e_r)\otimes e_s\bigr)\\
 &= (e_i\otimes e_j)
    \left(\sum_k\lambda_{kr}e_k\otimes e_s\right)\\
 &=\sum_k\lambda_{kr}(e_i\otimes e_j)(e_k\otimes e_s)\\
 &=\sum_k\lambda_{kr}\delta_{jk}\,e_i\otimes e_s\\
 &=\lambda_{jr}\,e_i\otimes e_s.
\end{align*}

Por consiguiente,
\begin{align*}
pTp
 &= (e_1\otimes e_1+e_2\otimes e_2)
    T
    (e_1\otimes e_1+e_2\otimes e_2)\\
 &= (e_1\otimes e_1)T(e_1\otimes e_1)
  +(e_1\otimes e_1)T(e_2\otimes e_2)\\
 &\qquad
  +(e_2\otimes e_2)T(e_1\otimes e_1)
  +(e_2\otimes e_2)T(e_2\otimes e_2)\\
 &=\lambda_{11}e_1\otimes e_1
  +\lambda_{12}e_1\otimes e_2
  +\lambda_{21}e_2\otimes e_1
  +\lambda_{22}e_2\otimes e_2.
\end{align*}

Así, todo elemento de $p\BH p$ queda determinado por cuatro coeficientes
complejos. Como espacio vectorial complejo,
\[
\dim_{\C}(p\BH p)=4.
\]
La base natural es
\[
\{e_1\otimes e_1,\ e_1\otimes e_2,\ e_2\otimes e_1,\ e_2\otimes e_2\}.
\]
Las reglas de producto
\[
(e_i\otimes e_j)(e_r\otimes e_s)=\delta_{jr}e_i\otimes e_s
\]
son las reglas de multiplicación matricial de $M_2(\C)$.

\begin{ejercicio}
Sea $S\in\BH$. Calcular explícitamente
\[
pSp.
\]
Después calcular
\[
(pSp)(pTp).
\]
\end{ejercicio}

\section{El ejemplo conmutativo $\C^n$}

En la parte final de la sesión se considera $\C^n$ con operaciones coordenada a
coordenada. Para
\[
z=(z_i)_{i=1}^n,
\qquad
w=(w_i)_{i=1}^n,
\]
se define
\[
zw=(z_iw_i)_{i=1}^n,
\qquad
z^*=(\overline{z_i})_{i=1}^n,
\]
y la norma
\[
\|z\|_\infty=\max_{1\le i\le n}|z_i|.
\]

Con estas operaciones, $\C^n$ es una $C^*$-álgebra conmutativa con unidad
\[
1=(1,\ldots,1).
\]
En las fotografías se relaciona este ejemplo con el teorema de representación
para $C^*$-álgebras conmutativas:
\[
\C^n\cong C(X),
\]
y se deja planteada la pregunta:
\begin{center}
\emph{¿Cuál es el espacio $X$ en este ejemplo?}
\end{center}

\section*{Notas de transcripción}

\begin{enumerate}[leftmargin=2.2em]
\item Las fotografías 9 y 10 contienen esencialmente el mismo desarrollo; se
transcribió una sola vez.

\item En el cálculo manuscrito de $p_iTp_i$ aparece en una de las fotografías un
coeficiente $\overline{\lambda_i}$. Con la convención escrita al inicio de la
misma sesión ---producto interno lineal en la primera variable--- y con
$T(e_i)=\sum_r\lambda_r e_r$, el cálculo directo da $\lambda_i$. Por ello, en el
texto principal se ha usado la versión consistente con la convención mostrada en
las fotografías, en lugar de modificarla silenciosamente.

\item En una anotación visual aparece $p\BH p\cong\C^4$ para la proyección de
rango dos. El cálculo de la clase sí muestra que el espacio tiene dimensión
compleja $4$. En estas notas se registra esa identificación como identificación
de espacios vectoriales; las reglas de producto escritas en la sesión son las de
las unidades matriciales y, por tanto, describen la estructura matricial
correspondiente.
\end{enumerate}

\end{document}
